\documentclass[11pt]{article}
\usepackage{reusv}
\usepackage{booktabs}
\usepackage{amsmath,amssymb}
\usepackage{siunitx}

\setborderthickness{0.5pt}
\setbordermargin{1cm}
\setbordercolor{black}

\slimtitle{아핀 드리프트 근사}

\begin{document}

\drawborder
\thispagestyle{firstpage}

\lightertitle{아핀 드리프트 근사:\\ 비선형 합성의 구간별 선형화}

{\normalsize\textit{Research Note No.~009 $|$ bezier-orbit-transfer-scp}}\par\medskip

\def\lighterdate{March 16, 2026}
\lighterauthor{Suwon Lee$^{\dagger}$}

\lighteraddress{$^\dagger$}{Future Mobility Control Lab, Kookmin University, Seoul, Korea}
\lighteremail{suwon.lee@kookmin.ac.kr}

\begin{abstract}
\cite{report006}의 SCP 정식화에서 중력 드리프트 적분 $\mathbf{c}_v$, $\mathbf{c}_r$은
참조 궤도를 수치 적분(RK4)하여 계산하며, 이것이 계산 시간의 86\%를 차지한다~\cite{report008}.
본 노트에서는 비선형 중력장 $\mathbf{a}_\text{grav}(\mathbf{r})$와
베지어 위치 궤적 $\mathbf{r}(\tau)$의 합성(composition)을
구간별 1차(또는 2차) 테일러 전개로 근사하여,
드리프트 적분을 제어점에 대한 아핀 연산으로 대체하는 방법을 유도한다.
드 카스텔조 세분화(de~Casteljau subdivision)를 통한
구간 분할이 자연스러운 trust-region 역할을 하며,
RK4 전파를 완전히 제거하여 SCP 반복당 계산 비용을
QP/SOCP 풀이 수준으로 줄일 수 있음을 보인다.
\end{abstract}

%% ============================================================
\section{문제 설정 및 동기}
\label{sec:motivation}
%% ============================================================

\subsection{현재 SCP 체계의 구조}

\cite{report006}에서 유도한 SCP 경계조건은 다음과 같다:
\begin{align}
  \mathbf{v}_f &= \mathbf{v}_0 + t_f \mathbf{c}_v + t_f \mathbf{I}_N \mathbf{P}_u
  \label{eq:bc-vel} \\
  \mathbf{r}_f &= \mathbf{r}_0 + t_f \mathbf{v}_0 + t_f^2 \mathbf{c}_r + t_f \bar{\mathbf{I}}_N \mathbf{P}_u
  \label{eq:bc-pos}
\end{align}
여기서 $\mathbf{I}_N$, $\bar{\mathbf{I}}_N$은 Bernstein 적분/이중적분 행렬~\cite{report004}이고,
드리프트 적분은
\begin{align}
  \mathbf{c}_v &= \int_0^1 \mathbf{a}_\text{grav}\bigl(\mathbf{r}_\text{ref}(\tau)\bigr)\,d\tau,
  \label{eq:cv-def} \\
  \mathbf{c}_r &= \int_0^1 \!\int_0^s \mathbf{a}_\text{grav}\bigl(\mathbf{r}_\text{ref}(\sigma)\bigr)\,d\sigma\,ds
  \label{eq:cr-def}
\end{align}
이다.

\subsection{병목 분석}

\cite{report008}의 프로파일링에 의하면, SCP 반복 1회당 계산 시간의 구성은
\begin{itemize}
  \item RK4 참조 궤도 전파: \textbf{86\%} (\SI{30}{\ms} / \SI{35}{\ms})
  \item CVXPY QP/SOCP 풀기: 11\% (\SI{4}{\ms})
  \item 기타: 3\%
\end{itemize}
이다. RK4가 지배적인 이유는, 식~\eqref{eq:cv-def}--\eqref{eq:cr-def}의
$\mathbf{a}_\text{grav}(\mathbf{r}_\text{ref}(\tau))$를 계산하려면
참조 궤도 $\mathbf{r}_\text{ref}(\tau)$를 전파해야 하기 때문이다.

\subsection{핵심 관찰}

위치 궤적 $\mathbf{r}(\tau)$는 초기조건과 적분 행렬을 통해
제어점 $\mathbf{P}_u$의 아핀 함수로 표현된다:
\begin{equation}
  \mathbf{r}(\tau) = \mathbf{r}_0 + t_f \mathbf{v}_0 \tau
    + t_f^2 \,\boldsymbol{\phi}(\tau)
    + t_f \, \mathbf{B}_{N+2}(\tau)^T \bar{\mathbf{I}}_N \, \mathbf{P}_u
  \label{eq:r-of-tau}
\end{equation}
여기서 $\boldsymbol{\phi}(\tau)$는 드리프트에 의한 이중 누적 적분(참조 궤도 의존)이다.
$\mathbf{r}(\tau)$가 $\mathbf{P}_u$에 대해 아핀이므로,
$\mathbf{a}_\text{grav}(\mathbf{r}(\tau))$를 $\mathbf{r}$ 주변에서
테일러 전개하면 $\mathbf{P}_u$에 대한 아핀(1차) 또는 이차(2차) 함수가 되며,
적분 역시 대수 연산으로 수행할 수 있다.

%% ============================================================
\section{비선형 합성의 테일러 근사}
\label{sec:taylor}
%% ============================================================

\subsection{중력장의 1차 전개}

참조 위치 $\mathbf{r}^*$ 주변에서:
\begin{equation}
  \mathbf{a}_\text{grav}(\mathbf{r}) \approx
    \mathbf{a}_\text{grav}(\mathbf{r}^*)
    + \mathbf{A}(\mathbf{r}^*)\,(\mathbf{r} - \mathbf{r}^*)
  \label{eq:taylor1}
\end{equation}
여기서 $\mathbf{A} = \partial \mathbf{a}_\text{grav}/\partial \mathbf{r} \in \mathbb{R}^{3\times 3}$은
중력 Jacobian이다.
2체 문제에서
\begin{equation}
  \mathbf{A}_\text{2B}(\mathbf{r}) = -\frac{1}{r^3}\mathbf{I}_3
    + \frac{3}{r^5}\mathbf{r}\mathbf{r}^T
\end{equation}
이며, J2 기여분은 \cite{report002}에서 유도한 해석적 Jacobian을 더한다.

\subsection{위치의 아핀 구조 대입}

식~\eqref{eq:r-of-tau}를 \eqref{eq:taylor1}에 대입하면,
참조 궤도에서의 위치 $\mathbf{r}^*(\tau) = \mathbf{r}_\text{ref}(\tau)$에 대해:
\begin{equation}
  \mathbf{a}_\text{grav}\bigl(\mathbf{r}(\tau)\bigr) \approx
    \underbrace{\mathbf{a}_\text{grav}\bigl(\mathbf{r}^*(\tau)\bigr)}_{\text{상수항 (참조 궤도)}}
    + \underbrace{\mathbf{A}\bigl(\mathbf{r}^*(\tau)\bigr) \cdot
      t_f \, \mathbf{B}_{N+2}(\tau)^T \bar{\mathbf{I}}_N}_{\text{$\mathbf{P}_u$에 대한 선형항}}
    \,\Delta \mathbf{P}_u
  \label{eq:agrav-affine}
\end{equation}
여기서 $\Delta\mathbf{P}_u = \mathbf{P}_u - \mathbf{P}_u^\text{ref}$이다.

\subsection{드리프트 적분의 아핀 근사}

식~\eqref{eq:agrav-affine}을 \eqref{eq:cv-def}에 대입하여 $\tau$에 대해 적분하면:
\begin{equation}
  \boxed{
  \mathbf{c}_v \approx \mathbf{c}_v^\text{ref}
    + \underbrace{t_f \int_0^1 \mathbf{A}\bigl(\mathbf{r}^*(\tau)\bigr)
      \mathbf{B}_{N+2}(\tau)^T \,d\tau \;\cdot\; \bar{\mathbf{I}}_N}_{\displaystyle \mathbf{M}_v \;\in\; \mathbb{R}^{3 \times (N+1)}}
    \;\Delta\mathbf{P}_u^{(k)}
  }
  \label{eq:cv-affine}
\end{equation}
여기서 $\mathbf{c}_v^\text{ref}$는 참조 궤도에서의 드리프트 적분(상수),
$\mathbf{M}_v$는 Jacobian--Bernstein 커플링 행렬이다.
마찬가지로
\begin{equation}
  \boxed{
  \mathbf{c}_r \approx \mathbf{c}_r^\text{ref}
    + \mathbf{M}_r \;\Delta\mathbf{P}_u^{(k)}
  }
  \label{eq:cr-affine}
\end{equation}
$\mathbf{M}_r$은 이중 적분에 대응하는 커플링 행렬이다.

\subsection{커플링 행렬의 수치 계산}

$\mathbf{M}_v$의 적분은 참조 궤도 위의 이산 점 $\{\tau_j\}_{j=0}^M$에서의
Gauss--Legendre 또는 사다리꼴 구적으로 계산한다:
\begin{equation}
  \mathbf{M}_v \approx t_f \sum_{j=0}^M w_j \,
    \mathbf{A}\bigl(\mathbf{r}^*(\tau_j)\bigr) \,
    \mathbf{B}_{N+2}(\tau_j)^T \;\bar{\mathbf{I}}_N
  \label{eq:Mv-quad}
\end{equation}
이 계산은 SCP 반복 시작 시 참조 궤도에서 \textbf{한 번만} 수행하며,
이후 QP의 경계조건 행렬에 반영된다.
$\mathbf{A}(\mathbf{r}^*(\tau_j))$는 해석적 Jacobian이므로 빠르게 계산 가능하다.

\subsection{2차 전개 (선택적)}

더 높은 근사 정확도가 필요할 경우, 2차 항을 포함할 수 있다:
\begin{equation}
  \mathbf{a}_\text{grav}(\mathbf{r}) \approx
    \mathbf{a}_\text{grav}(\mathbf{r}^*)
    + \mathbf{A}(\mathbf{r}^*)\delta\mathbf{r}
    + \tfrac{1}{2}\,\mathbf{H}(\mathbf{r}^*)\,[\delta\mathbf{r}, \delta\mathbf{r}]
\end{equation}
여기서 $\mathbf{H} = \partial^2 \mathbf{a}_\text{grav}/\partial \mathbf{r}^2$는
중력 Hessian 텐서이다.
$\delta\mathbf{r}$이 $\mathbf{P}_u$에 아핀이므로,
2차 항은 $\Delta\mathbf{P}_u$에 대한 이차형식이 되어 QCQP로 정식화 가능하다.
다만, 구간 세분화로 1차 근사의 정확도를 높이는 것이 일반적으로 더 효율적이다.

%% ============================================================
\section{드 카스텔조 세분화}
\label{sec:subdivision}
%% ============================================================

\subsection{세분화와 구간별 선형화}

$N$차 베지어 곡선을 드 카스텔조 알고리즘으로 $\tau = \tau_k$에서 분할하면,
원래 곡선과 정확히 동일한 두 개의 $N$차 베지어 곡선이 생성된다.
이를 반복하면 $[0,1]$을 $K$개 구간 $[\tau_{k-1}, \tau_k]$으로 분할하고,
각 구간을 별도의 $N$차 베지어 곡선 $\mathbf{q}_k(s)$, $s\in[0,1]$로 표현한다.

각 구간 $k$에서의 참조점 $\mathbf{r}_k^*$(예: 구간 중점에서의 참조 위치)를 잡으면:
\begin{equation}
  \mathbf{a}_\text{grav}\bigl(\mathbf{r}_k(s)\bigr) \approx
    \mathbf{a}_\text{grav}(\mathbf{r}_k^*)
    + \mathbf{A}(\mathbf{r}_k^*)\,\bigl(\mathbf{r}_k(s) - \mathbf{r}_k^*\bigr)
  \label{eq:piecewise-taylor}
\end{equation}

\subsection{근사 오차 제어}

전체 구간에서 하나의 참조점을 사용하면 테일러 잔여항이
$O(\|\mathbf{r} - \mathbf{r}^*\|^2)$이다.
$K$개 구간으로 세분화하면 각 구간의 잔여항이
$O(\|\mathbf{r}_k - \mathbf{r}_k^*\|^2)$인데,
구간이 짧을수록 $\|\mathbf{r}_k - \mathbf{r}_k^*\|$가 작으므로
오차가 $O(1/K^2)$로 감소한다 (1차 근사 기준).

이는 SCP에서 trust-region을 외부 제약으로 부과하는 대신,
\textbf{세분화 수준 $K$가 자연스러운 trust-region 반경을 결정}함을 의미한다:
\begin{itemize}
  \item $K$ 증가 $\to$ 각 구간의 선형화 유효 범위 확대 $\to$ trust-region 확장
  \item $K$ 감소 $\to$ 계산량 감소, 단 선형화 오차 증가
\end{itemize}

\subsection{세분화된 드리프트 적분}

전체 드리프트 적분을 구간별 합산으로 표현한다:
\begin{equation}
  \mathbf{c}_v = \sum_{k=1}^K \Delta\tau_k \int_0^1
    \mathbf{a}_\text{grav}\bigl(\mathbf{r}_k(s)\bigr)\,ds
  \approx \sum_{k=1}^K \Delta\tau_k \Bigl[
    \mathbf{a}_\text{grav}(\mathbf{r}_k^*)
    + \mathbf{A}(\mathbf{r}_k^*)\,\bigl(\overline{\mathbf{r}}_k - \mathbf{r}_k^*\bigr)
  \Bigr]
  \label{eq:cv-piecewise}
\end{equation}
여기서 $\overline{\mathbf{r}}_k = \int_0^1 \mathbf{r}_k(s)\,ds$는
베지어 제어점의 산술평균으로 \textbf{해석적}으로 계산된다~\cite{report004}:
\begin{equation}
  \overline{\mathbf{r}}_k = \frac{1}{N+1}\sum_{i=0}^N \mathbf{p}_{k,i}
\end{equation}

따라서 구간별 드리프트 적분이 완전히 대수적 연산이 된다.

%% ============================================================
\section{수정된 SCP 구조}
\label{sec:modified-scp}
%% ============================================================

\subsection{알고리즘 개요}

아핀 드리프트 근사를 적용한 SCP 반복은 다음과 같다:

\begin{enumerate}
  \item \textbf{참조 궤도 평가} (반복 시작 시 1회):
    \begin{itemize}
      \item 현재 제어점 $\mathbf{P}_u^\text{ref}$로부터 이산 점
            $\{\mathbf{r}^*(\tau_j)\}$를 계산
            (식~\eqref{eq:r-of-tau}의 대수적 평가, RK4 불필요)
      \item 각 이산 점에서 $\mathbf{a}_\text{grav}(\mathbf{r}^*)$와
            Jacobian $\mathbf{A}(\mathbf{r}^*)$ 계산
    \end{itemize}
  \item \textbf{커플링 행렬 구성}: $\mathbf{M}_v$, $\mathbf{M}_r$ 계산
        (식~\eqref{eq:Mv-quad}, 행렬 연산)
  \item \textbf{경계조건 행렬 수정}: $\mathbf{c}_v$, $\mathbf{c}_r$의 선형 기여분을
        $\mathbf{A}_\text{eq}$에 흡수
  \item \textbf{QP/SOCP 풀기}: 수정된 경계조건으로 볼록 부분문제 풀기
  \item \textbf{수렴 판정}: $\|\Delta\mathbf{P}_u\|$ 및 실제 BC 위반 확인
\end{enumerate}

\subsection{경계조건 행렬 수정}

식~\eqref{eq:cv-affine}을 \eqref{eq:bc-vel}에 대입하면:
\begin{equation}
  \mathbf{v}_f = \mathbf{v}_0 + t_f\mathbf{c}_v^\text{ref}
    + t_f(\mathbf{I}_N + \mathbf{M}_v)\,\mathbf{P}_u
    - t_f \mathbf{M}_v\,\mathbf{P}_u^\text{ref}
\end{equation}
즉, 기존 적분 행렬 $\mathbf{I}_N$에 커플링 행렬 $\mathbf{M}_v$가 더해지고,
상수항에 참조 제어점 기여분이 추가된다.
보조변수 $\mathbf{Z} = t_f\mathbf{P}_u$를 사용하면:
\begin{equation}
  (\mathbf{I}_N + \mathbf{M}_v)\,\mathbf{Z} =
    (\mathbf{v}_f - \mathbf{v}_0) - t_f\mathbf{c}_v^\text{ref}
    + \mathbf{M}_v\,\mathbf{Z}^\text{ref}
  \label{eq:bc-vel-modified}
\end{equation}
이는 여전히 $\mathbf{Z}$에 대한 \textbf{선형 등식 제약}이다.
위치 제약도 동일한 구조로 수정된다.

\subsection{계산 비용 비교}

\begin{table}[h]
\centering
\caption{SCP 반복 1회당 계산 비용 비교.}
\label{tab:cost-compare}
\small
\begin{tabular}{@{}lcc@{}}
\toprule
단계 & 현재 (0차 + RK4) & 제안 (1차 아핀) \\
\midrule
참조 궤도 평가 & RK4 전파 (\SI{15}{\ms}) & 행렬 곱 ($<\SI{0.1}{\ms}$) \\
드리프트 적분   & 수치 적분 (\SI{1}{\ms}) & $\mathbf{M}_v$, $\mathbf{M}_r$ 구성 (\SI{0.5}{\ms}) \\
QP/SOCP        & CVXPY (\SI{4}{\ms})      & CVXPY (\SI{4}{\ms}) \\
BC 검증        & RK4 재전파 (\SI{15}{\ms}) & 대수적 평가 ($<\SI{0.1}{\ms}$) \\
\midrule
합계           & $\sim$\SI{35}{\ms}       & $\sim$\SI{5}{\ms} \\
\bottomrule
\end{tabular}
\end{table}

RK4 전파가 완전히 제거되어 \textbf{약 7배} 속도 향상이 기대된다.

%% ============================================================
\section{구현 방향}
\label{sec:implementation}
%% ============================================================

\subsection{Phase 1: 1차 아핀 드리프트 (세분화 없음)}

전체 구간에서 단일 참조점을 사용하는 가장 단순한 버전.
기존 SCP 루프에서 RK4 전파를 대수적 평가로 교체한다.

\begin{enumerate}
  \item \texttt{bezier/basis.py}: 이산 점 $\tau_j$에서의 $\mathbf{B}_{N+2}(\tau_j)$ 일괄 평가 함수 추가
  \item \texttt{scp/linearize.py} (신규): 커플링 행렬 $\mathbf{M}_v$, $\mathbf{M}_r$ 계산
  \item \texttt{scp/inner\_loop.py}: \texttt{\_compute\_drift\_integrals}를 아핀 버전으로 교체
  \item \texttt{scp/problem.py}: 수정된 경계조건 행렬 빌더
\end{enumerate}

검증: \cite{report008}의 V1--V3 시나리오에서 동일한 해를 산출하는지 확인.

\subsection{Phase 2: 드 카스텔조 세분화}

\begin{enumerate}
  \item \texttt{bezier/subdivision.py} (신규): 드 카스텔조 분할 알고리즘
    \begin{itemize}
      \item \texttt{de\_casteljau\_split(P, tau)}: 제어점 $\mathbf{P}$를 $\tau$에서 분할
      \item \texttt{subdivide\_uniform(P, K)}: $K$개 균등 구간 분할
      \item \texttt{subdivide\_adaptive(P, tol)}: 곡률 기반 적응형 분할
    \end{itemize}
  \item \texttt{scp/linearize.py} 확장: 구간별 $\mathbf{M}_{v,k}$, $\mathbf{M}_{r,k}$ 계산 및 합산
  \item Trust-region 제약 제거: 세분화가 자연 trust-region 역할
\end{enumerate}

\subsection{Phase 3: 적응형 세분화 및 수렴 가속}

\begin{enumerate}
  \item 적응형 $K$ 결정: 선형화 잔여항 추정으로 세분화 수준 자동 조절
  \item 수렴 특성 분석: 0차 vs.\ 1차 + 세분화의 반복 횟수 비교
  \item 수렴 판정 개선: 아핀 BC 잔차를 직접 사용 (전파 불필요)
\end{enumerate}

\subsection{핵심 데이터 구조}

\begin{verbatim}
@dataclass
class AffineLinearization:
    """구간별 1차 선형화 결과."""
    K: int                    # 세분화 구간 수
    tau_breaks: NDArray       # (K+1,) 구간 경계
    a_ref: list[NDArray]      # K × (3,) 참조점 중력 가속도
    A_ref: list[NDArray]      # K × (3,3) 참조점 Jacobian
    r_ref: list[NDArray]      # K × (3,) 참조점 위치
    Mv: NDArray               # (3, N+1) 속도 커플링 행렬
    Mr: NDArray               # (3, N+1) 위치 커플링 행렬
    cv_ref: NDArray           # (3,) 참조 드리프트 적분 (속도)
    cr_ref: NDArray           # (3,) 참조 드리프트 적분 (위치)
\end{verbatim}

%% ============================================================
\section{수치 예제}
\label{sec:numerical}
%% ============================================================

구간별 테일러 근사가 RK4 기반 드리프트 적분을 얼마나 정확히
재현하는지 검증한다.
수렴된 SCP 해의 참조 궤도 위에서, 중력 가속도 $\mathbf{a}_\text{grav}(\mathbf{r}(\tau))$의
적분을 $K$개 구간의 구간 중점 테일러 전개로 근사하여
RK4 결과(201스텝 사다리꼴 구적)와 비교한다.

\subsection{시나리오 및 설정}

세 가지 시나리오를 사용한다:

\begin{itemize}
  \item \textbf{A}: LEO $\to 1.1a_0$ 원형 (약한 전이, $t_f=3.0$)
  \item \textbf{B}: LEO $\to 1.2a_0$ 원형 (중간 전이, $t_f=3.62$)
  \item \textbf{C}: LEO $\to$ LEO, $\Delta i=10^\circ$ (경사각 변경, $t_f=3.0$)
\end{itemize}

모두 $N=12$, J2 섭동 포함, 정규화 단위 기준이다.
비교 대상은 두 가지 근사 방식이다:
\begin{itemize}
  \item \textbf{0차}: 각 구간 중점의 $\mathbf{a}_\text{grav}(\mathbf{r}_k^*)$ 값만 사용
        (구간 midpoint rule)
  \item \textbf{1차}: 중점 Jacobian 보정 포함:
        $\mathbf{a}_\text{grav}(\mathbf{r}_k^*) + \mathbf{A}(\mathbf{r}_k^*)\,(\mathbf{r} - \mathbf{r}_k^*)$
\end{itemize}

\subsection{$\mathbf{c}_v$ 근사 정확도}

\begin{table}[h]
\centering
\caption{$\mathbf{c}_v$ 상대 오차 (\%): $K$개 구간 근사 vs.\ RK4 기준.}
\label{tab:cv-accuracy}
\small
\begin{tabular}{@{}c cc cc cc@{}}
\toprule
 & \multicolumn{2}{c}{시나리오 A} & \multicolumn{2}{c}{시나리오 B} & \multicolumn{2}{c}{시나리오 C} \\
\cmidrule(lr){2-3} \cmidrule(lr){4-5} \cmidrule(lr){6-7}
$K$ & 0차 & 1차 & 0차 & 1차 & 0차 & 1차 \\
\midrule
1  & 5.34 & 6.21 & 6.77 & 8.12 & 6.74 & 7.52 \\
2  & 4.92 & 1.88 & 6.65 & 2.25 & 5.48 & 2.25 \\
4  & 0.52 & 1.00 & 0.88 & 1.28 & 0.52 & 1.17 \\
8  & 0.073 & 0.31 & 0.16 & 0.41 & 0.048 & 0.36 \\
16 & 0.017 & 0.081 & 0.037 & 0.11 & 0.009 & 0.094 \\
32 & 0.004 & 0.021 & 0.009 & 0.028 & 0.003 & 0.025 \\
\bottomrule
\end{tabular}
\end{table}

\subsection{$\mathbf{c}_r$ 근사 정확도}

\begin{table}[h]
\centering
\caption{$\mathbf{c}_r$ 상대 오차 (\%): $K$개 구간 근사 vs.\ RK4 기준.}
\label{tab:cr-accuracy}
\small
\begin{tabular}{@{}c cc cc cc@{}}
\toprule
 & \multicolumn{2}{c}{시나리오 A} & \multicolumn{2}{c}{시나리오 B} & \multicolumn{2}{c}{시나리오 C} \\
\cmidrule(lr){2-3} \cmidrule(lr){4-5} \cmidrule(lr){6-7}
$K$ & 0차 & 1차 & 0차 & 1차 & 0차 & 1차 \\
\midrule
1  & 10.5 & 6.05 & 12.6 & 7.70 & 12.3 & 7.52 \\
2  & 7.00 & 1.94 & 9.55 & 2.40 & 6.58 & 2.33 \\
4  & 1.56 & 1.11 & 2.22 & 1.52 & 1.36 & 1.23 \\
8  & 0.38 & 0.34 & 0.53 & 0.48 & 0.33 & 0.36 \\
16 & 0.10 & 0.088 & 0.15 & 0.12 & 0.082 & 0.095 \\
32 & 0.034 & 0.023 & 0.055 & 0.033 & 0.023 & 0.025 \\
\bottomrule
\end{tabular}
\end{table}

\subsection{분석}

\begin{enumerate}
  \item \textbf{0차 근사의 수렴}: 구간 midpoint rule에 해당하며,
        매끄러운 피적분 함수에 대해 $O(1/K^2)$ 수렴을 보인다.
        수치적으로 $\approx O(1/K^{2.1\text{--}2.5})$이 관찰된다.
  \item \textbf{1차 근사}: 동일한 고정 참조 궤도 위에서의 구적(quadrature)
        정확도에서는 0차와 비슷하거나 약간 큰 오차를 보인다.
        이는 midpoint rule의 최적 오차 상쇄 특성 때문이다.
        1차의 진정한 장점은 제어점 $\mathbf{P}_u$ 변화에 대한
        감도(sensitivity)를 반영하는 데 있으며,
        이는 제어점 업데이트 시 드리프트를 재계산 없이 보정할 수 있게 한다.
  \item \textbf{실용적 권장}: $K=8$이면 $\mathbf{c}_v$, $\mathbf{c}_r$ 모두
        \textbf{0.5\% 이내} 오차로 RK4 결과를 재현한다.
        이는 200스텝 RK4 전파를 8개 참조점의 중력 평가로 대체할 수 있음을 의미하며,
        중력 가속도 및 Jacobian 평가 비용이 $O(K)$이므로 대폭적인 계산량 절감이 가능하다.
  \item \textbf{이중 적분 $\mathbf{c}_r$}: 단일 적분 $\mathbf{c}_v$보다 오차가 크지만,
        1차 보정 시 $K=4$에서도 약 1\% 수준으로 충분히 정확하다.
\end{enumerate}

\subsection{계산 시간 비교}

시나리오 B에서 드리프트 적분 계산 시간을 측정한다 (50회 반복 평균).

\begin{table}[h]
\centering
\caption{드리프트 적분 계산 시간 비교 (시나리오 B, $N=12$).}
\label{tab:timing}
\small
\begin{tabular}{@{}lcccc@{}}
\toprule
방법 & 시간 (ms) & 속도 향상 & $\mathbf{c}_v$ 오차 & $\mathbf{c}_r$ 오차 \\
\midrule
RK4 전파 + 수치 적분 & 15.35 & 기준 & --- & --- \\
\quad $\llcorner$ RK4 전파 (201 steps)  & 15.31 & & & \\
\quad $\llcorner$ 드리프트 적분 (벡터화) & 0.04 & & & \\
\midrule
0차 midpoint $K=4$ & 0.042 & $365\times$ & 0.88\% & 4.4\% \\
0차 midpoint $K=8$ & 0.042 & $370\times$ & 0.16\% & 1.2\% \\
0차 midpoint $K=16$ & 0.043 & $347\times$ & 0.037\% & 0.29\% \\
0차 midpoint $K=32$ & 0.044 & $348\times$ & 0.009\% & 0.071\% \\
\midrule
1차 Taylor $K=4$ & 0.16 & $96\times$ & 1.3\% & --- \\
1차 Taylor $K=8$ & 0.31 & $50\times$ & 0.41\% & --- \\
1차 Taylor $K=16$ & 0.60 & $26\times$ & 0.11\% & --- \\
\bottomrule
\end{tabular}
\end{table}

주목할 결과:
\begin{enumerate}
  \item \textbf{RK4의 병목은 전파 자체}: 15.35\,ms 중 15.31\,ms가 RK4 전파이며,
        벡터화된 드리프트 적분은 0.04\,ms에 불과하다.
        즉, RK4 전파를 제거하면 드리프트 적분 비용은 무시할 수 있다.
  \item \textbf{0차 midpoint의 효율}: $K$에 거의 무관하게 $\sim$0.04\,ms로,
        RK4 대비 \textbf{약 370배} 빠르다.
        벡터화된 중력 평가($K$점 일괄 처리)가 매우 효율적이기 때문이다.
  \item \textbf{1차 Taylor의 비용}: Jacobian 계산($3\times 3$ 행렬 $\times K$회)과
        구적점 보간 때문에 $K$에 선형 비례한다.
        $K=8$에서도 0.31\,ms로 RK4 대비 50배 빠르지만,
        0차 대비 7배 느리다.
  \item \textbf{SCP 반복당 총 효과}: 현재 반복당 $\sim$19\,ms 중 RK4 전파가
        $\sim$15\,ms를 차지하므로,
        이를 0차 midpoint로 대체하면 반복당 $\sim$4\,ms로 단축된다.
\end{enumerate}

%% ============================================================
\section{이론적 고찰}
\label{sec:theory}
%% ============================================================

\subsection{SCP 수렴 차수와의 관계}

현재 0차 선형화(드리프트를 상수로 취급)에서는 SCP 반복이 고정점 반복(fixed-point iteration)의
성격을 가지며, 수렴 속도가 선형이다~\cite{report008}.
1차 아핀 근사를 도입하면 Newton-type 반복에 가까워져
초선형(superlinear) 수렴이 기대된다.

\subsection{볼록성 보존}

식~\eqref{eq:bc-vel-modified}에서 볼 수 있듯이, 아핀 드리프트 근사 후에도
경계조건은 $\mathbf{Z}$에 대한 선형 등식 제약을 유지한다.
목적함수 $J = (1/t_f)\sum_k \mathbf{z}_k^T\mathbf{G}_N\mathbf{z}_k$와
추력 SOCP 제약은 변경되지 않으므로,
부분문제의 볼록성이 완전히 보존된다.

\subsection{세분화와 정확도-효율 트레이드오프}

$K$개 구간 세분화 시:
\begin{itemize}
  \item 선형화 오차: $O(1/K^2)$ (1차), $O(1/K^3)$ (2차)
  \item 커플링 행렬 계산 비용: $O(K \cdot M \cdot (N+1))$
        (각 구간 $M$개 구적점, $(N+1)$개 기저 평가)
  \item SCP 반복 횟수: $K$ 증가에 따라 감소 (더 정확한 선형화)
\end{itemize}

실용적으로 $K = 4$--$8$ 정도면 RK4 200스텝 전파 대비
충분한 정확도를 확보할 것으로 예상된다.

%% ============================================================
\section{결론}
\label{sec:conclusion}
%% ============================================================

베지어 곡선으로 파라미터화된 위치 궤적 $\mathbf{r}(\tau)$가 제어점의 아핀 함수라는
구조적 성질을 활용하여, 비선형 중력장과의 합성을 구간별 테일러 전개로 근사하는
아핀 드리프트 기법을 유도하였다.

\begin{enumerate}
  \item 드리프트 적분 $\mathbf{c}_v$, $\mathbf{c}_r$이 제어점에 대한
        아핀 함수로 근사되어, RK4 수치 적분이 행렬 연산으로 대체된다.
  \item 드 카스텔조 세분화가 구간별 선형화의 정확도를 제어하며,
        외부 trust-region 제약의 자연스러운 대안이 된다.
  \item SCP 부분문제의 볼록성이 완전히 보존되며,
        경계조건 행렬에 커플링 항 $\mathbf{M}_v$, $\mathbf{M}_r$만 추가된다.
\end{enumerate}

Phase~1 구현(단일 구간 1차 근사)부터 시작하여, V1--V3 검증 시나리오에서의
정합성을 확인한 후, Phase~2 드 카스텔조 세분화로 확장하는 단계적 전략을 취한다.

\bibliographystyle{IEEEtran}
\bibliography{references}

\end{document}
